\documentclass[10pt,a4paper]{amsart}
\usepackage[margin=19mm]{geometry}
\usepackage{amsmath,amssymb,mathtools}
\usepackage[scr=rsfs,cal=euler]{mathalfa}
\usepackage{xfrac}
\usepackage[unicode,hidelinks,bookmarks=false]{hyperref}
\newcommand{\ie}{i.e.\ }
\newcommand{\cf}{cf.\ }
\newcommand{\resp}{respectively\ }
\newcommand{\Subsection}{\subsection}
\providecommand{\nobreakdash}{\nobreak\hbox{-}}
\setlength{\emergencystretch}{2em}
\multlinegap0pt
\usepackage{bm}
\usepackage{bbm}
\usepackage{dsfont}
\usepackage{xspace}
\usepackage{cancel}
\usepackage{mathtools}
\usepackage{amsthm}
\makeatletter
\@ifundefined{theorem}{\newtheorem{theorem}{Theorem}}{}
\makeatother
\makeatletter
\def\DIS{\mathsf{DIS}}
\def\IN{{\Bbb{N}}}
\def\In{\subseteq}
\expandafter\ifx\csname cases\endcsname\relax
\newenvironment{cases}{\left\{\begin{array}{ll}}{\end{array}\right.}
\fi
\def\dom{{\mathrm{dom}}}
\expandafter\ifx\csname eqcase\endcsname\relax
\newenvironment{eqcase}{\left\{\begin{array}{lcl}}{\end{array}\right.}
\fi
\def\equivSW{\mathop{\equiv_{\mathrm{sW}}}}
\def\graph{{\mathrm{graph}}}
\def\id{{\mathrm{id}}}
\def\mto{\rightrightarrows}
\def\prefix{\sqsubseteq}
\def\r{\mathrm{r}}
\def\range{{\mathrm{range}}}
\makeatother
\title[The Discontinuity Problem]{The Discontinuity Problem (Theorem 6)}
\author{Vasco Brattka}
\date{Source: arXiv:2012.02143v1 (CC BY 4.0)}
\begin{document}
\maketitle

\noindent\emph{Source and licence.} The Discontinuity Problem. Version arXiv:2012.02143v1 (CC BY 4.0), distributed under the Creative Commons Attribution 4.0 International licence, as recorded in the arXiv licence metadata for this exact version and shown on the abstract page https://arxiv.org/abs/2012.02143v1. Full rights notice: licenses/arxiv-2012.02143-CC-BY-4.0.txt.\par\smallskip

\noindent\textit{Editorial scope.} Counted unit: the Theorem printed as Theorem 6 of the article (source0.tex lines 1116--1124, source label \texttt{thm:Wadge-game}), delivered together with the article's own complete original proof of it (source0.tex lines 1125--1166, one \texttt{proof} environment of 42 lines). No mathematics is written, repaired or re-derived: statement and proof are the article's own text line for line. Required context retained uncounted: thm:dichotomy (lines 557--564). Every definition, display and label of those lines is retained unchanged. Omitted from this package and not redistributed: 1--556, 565--1115, 1167--1551. Nothing inside a retained excerpt is omitted or altered: every retained body is delivered exactly as the article's lines are written, so no omission receipt of any kind accompanies this package. Citations: the retained excerpts carry no citation command at all, so no bibliography entry of the article is transcribed and no reference is redistributed. Reference adaptation: none was needed and none was made; every reference inside the delivered unit resolves inside it, so no unresolved reference or citation is produced. Printed slips kept literal: no printed word, symbol, hypothesis or number of the counted statement or of its proof was altered. Layout and class: the article's own class is \texttt{amsart} with packages including hyperref, amsmath, amssymb, amsthm, tikz, catchfile, lineno; the wrapper is the lane's pinned \texttt{amsart} editorial wrapper at 10pt on A4, which restates the theorem-like environments the retained text needs and adds the article's own macro definitions that the retained text needs (\texttt{\textbackslash DIS}, \texttt{\textbackslash IN}, \texttt{\textbackslash In}, \texttt{\textbackslash dom}, \texttt{\textbackslash equivSW}, \texttt{\textbackslash graph}, \texttt{\textbackslash id}, \texttt{\textbackslash mto}, \texttt{\textbackslash prefix}, \texttt{\textbackslash r}, \texttt{\textbackslash range}), with extra wrapper packages (bm, bbm, dsfont, xspace, cancel, mathtools). The delivered numbering is the unit's own. Assets: no retained span contains a figure environment, a graphics-inclusion command, a PDF-page inclusion, a table or a TikZ picture; no image file is copied into this package, the package declares no asset dependency and no image file is redistributed with it. Licence: version arXiv:2012.02143v1 of this article is distributed under the Creative Commons Attribution 4.0 International licence, as recorded in the arXiv licence metadata for this exact version and shown on the abstract page https://arxiv.org/abs/2012.02143v1. A full rights notice is delivered at licenses/arxiv-2012.02143-CC-BY-4.0.txt. Characters: every retained excerpt is pure ASCII, so the delivered engine, XeLaTeX with its default Latin Modern text font, reports no missing character.
\begin{theorem}[Continuity and effective discontinuity]
\label{thm:dichotomy}
Let $f:\In X\mto Y$ be a problem. Then we obtain
\begin{enumerate}
\item $f\leq_\mathrm{W}^*\id\iff f$ is continuous,
\item $\DIS\leq_\mathrm{W}^*f\iff f$ is effectively discontinuous.
\end{enumerate}
\end{theorem}

\begin{theorem}[Wadge games]
\label{thm:Wadge-game}
We consider the Wadge game of a given problem ${f:\In X\mto Y}$.
Then the following hold:
\begin{enumerate}
\item $f$ is continuous $\iff$ Player II has a winning strategy for $f$,
\item $f$ is effectively discontinuous $\iff$ Player I has a winning strategy for $f$.
\end{enumerate}
\end{theorem}

\begin{proof}
Since the Wadge game of $f$ is the Wadge game of $f^\r$ and $f^\r\equivSW f$, it suffices by 
Theorem~\ref{thm:dichotomy} to consider problems of type $f:\In\IN^\IN\mto\IN^\IN$.\smallskip\\
(1) If $f$ is continuous, then $f$ has a continuous realizer $F:\In\IN^\IN\to\IN^\IN$, which is approximated
by a monotone function $h:\IN^*\to\IN^*$ in the sense that $F(p)=\sup_{w\prefix p}h(w)$ for all $p\in\dom(F)\supseteq\dom(f)$. Given the moves $w_0,w_1,...\in\IN^*$ of Player~I, we can
inductively define the moves $v_i\in\IN^*$ by $v_0...v_i:=h(w_0...w_i)$ for all $i\in\IN$, since $h$ is monotone.
Then $\sigma(\varepsilon):=\varepsilon$ and $\sigma(w_0,...,w_i):=v_i$ provides a winning strategy
$\sigma:\IN^{**}\to\IN^*$ for Player~II.
This is because if $r:=w_0w_1...\in\dom(f)$, then $F(r)=v_0v_1...\in f(r)$.\smallskip\\
Vice versa, let $\sigma:\IN^{**}\to\IN^*$ be a winning strategy for Player~II.
Then we can define $h:\IN^*\to\IN^*$ by $h(\varepsilon):=\varepsilon$ and $h(a_0...a_i):=v_0...v_i$ for all $a_0,a_1,...\in\IN$,
where we inductively choose  $v_i:=\sigma(a_0,...,a_i)$.
Let $F:\In\IN^\IN\to\IN^\IN$ be given by $F(p):=\sup_{w\prefix p}h(w)$.
Given an input $r:=a_0a_1...\in\dom(f)$, we obtain $F(r)=v_0v_1...$ such that
$(r,F(r))\in\graph(f)$, since $\sigma$ is a winning strategy for Player II. 
Hence $F$ is a continuous realizer for $f$.\medskip\\
(2) If $f$ is effectively discontinuous, then there is a continuous $D:\IN^\IN\to\IN^\IN$
that witnesses the discontinuity of $f$ in the sense that $D(p)\in\dom(f)$ and
$\Phi_pD(p)\not\in fD(p)$ for all $p\in\IN^\IN$.
Let $h:\IN^*\to\IN^*$ be a monotone function that approximates $D$, in the sense
that $D(p)=\sup_{w\prefix p}h(w)$ for all $p\in\IN^\IN$.
Given the moves $v_0,v_1,...\in\IN^*$ of Player II, we can inductively
define the moves $w_i$ by $w_0:=h(\varepsilon)$ and  $w_0...w_i:=h(\langle\overline{w_0},\overline{v_0}\rangle...\langle\overline{w_0...w_{i-1}},\overline{v_0...v_{i-1}}\rangle)$. 
We note that $r:=w_0w_1...$ is an infinite sequence $r\in\dom(f)$,
since $D$ is total and $\range(D)\In\dom(f)$. And $p:=\langle\overline{w_0},\overline{v_0}\rangle\langle\overline{w_0w_{1}},\overline{v_0v_{1}}\rangle...\in\IN^\IN$
is the name of a function $\Phi_p$ with $D(p)=r$. 
If $q:=v_0v_1...$ is finite, then clearly $(r,q)\not\in\graph(f)$. Otherwise, $q=\Phi_p(r)=\Phi_pD(p)\not\in fD(p)=f(r)$
and hence also $(r,q)\not\in\graph(f)$. This means that $\sigma:\IN^{**}\to\IN^*$
with $\sigma(v_0,...,v_{i-1}):=w_i$ is a winning strategy for Player~I.\smallskip\\
Vice versa let  $\sigma:\IN^{**}\to\IN^*$ be a winning strategy for Player~I.
We need to define a continuous discontinuity function $D:\IN^\IN\to\IN^\IN$ for $f$.
Given $p\in\IN^\IN$, we can determine a monotone $h:\IN^*\to\IN^*$ that approximates $\Phi_p$, i.e., 
such that $\Phi_p(q)=\sup_{w\prefix q}h(w)$ for all $q\in\dom(\Phi_p)$. 
We let $D(p):=w_0w_1...$
where the $w_i$ are inductively given by $w_i:=\sigma(v_0,...,v_{i-1})$
and $v_0...v_i:=h(w_0...w_i)$.
Since $\sigma$ is a winning strategy for Player~I, we have $D(p)\in\dom(f)$. In particular, $D$ is total and continuous.
Moreover, with $q:=v_0v_1...$ we have $(D(p),q)\not\in\graph(f)$.
This could mean that $q$ is finite and hence $D(p)\not\in\dom(\Phi_p)$ 
or otherwise $\Phi_pD(p)=q$. In any case, $\Phi_pD(p)\not\in fD(p)$ holds and $D$ is a discontinuity function
for $f$.
\end{proof}

\end{document}
